
\section{Введение}
По аналогии с поверхностью Маркова, мы вводим более общее понятие поверхностей марковского типа и находим их группы автоморфизмов. 
Раздел 2 содержит предварительные сведения о пополнениях аффинных поверхностей, их бирациональных преобразованиях и двойственных графах граничных дивизоров. Раздел 3 содержит описание пополнения, граничного дивизора, связи с группой Томпсона и группы автоморфизмов для каждой из рассматриваемых поверхностей.

%{\em Благодарность.} Автор благодарен Владу Сержиеску за доступное введение в группу Томпсона. Автор особенно благодарен Михаилу Григорьевичу Зайденбергу за погружение в тематику автоморфизмов аффинных поверхностей, а также за многочисленные обсуждения, замечания и предложения по статье. Наконец, автор благодарен Юрию Геннадьевичу Прохорову за замечания, обогатившие статью.

\section{Предварительные сведения о пополнениях и двойственных графах}


\subsection{Пополнения аффинных поверхностей}
Пусть $\K$ --- алгебраически замкнутое поле характеристики 0, и $Y$ --- неприводимая аффинная алгебраическая поверхность над $\K$. \emph{Пополнением} $Y$ будем называть открытое вложение $Y\embed X$ в неприводимую проективную поверхность $X$. Например, если $Y\subset \AA^n\subset\PP^n$, то можем взять $X=\overline{Y}$ --- замыкание в топологии Зарисского в проективном пространстве $\PP^n$. 

Мы предполагаем, что \emph{граничный дивизор} пополнения $D=X\setminus Y$ имеет \emph{простые нормальные пересечения}, то есть каждая особая точка дивизора является трансверсальным пересечением пары его различных неприводимых компонент (то есть кривых). При этих условиях пара $(X,D)$ называется \emph{NC-парой}.  Дивизору $D$ можно сопоставить взвешенный \emph{двойственный граф} $\Gamma(D)$, взяв
\begin{itemize}
    \item в качестве вершин $\Gamma(D)$ --- неприводимые компоненты $D$;
    \item в качестве весов вершин --- индексы самопересечения компонент;
    \item в качестве рёбер --- точки пересечения компонент $D$.
\end{itemize}
Заметим, что двойственный граф может иметь кратные рёбра.

Автоморфизмы аффинной поверхности $Y$ продолжаются как бирациональные преобразования пополнения $X$ и описываются в терминах раздутий и стягиваний граничного дивизора $D$. 
Группа бирациональных преобразований $X$ в себя, регулярных на $Y$, обозначается через $\Bir(X,D)$. Она изоморфна $\Aut(Y)$.
Индуцированные \emph{бирациональные} преобразования двойственного графа $\Gamma(D)$ описаны в \cite{FKZ07a}.

\subsection{Взвешенные графы}

Мы напомним необходимые сведения и результаты о двойственных графах из \cite{FKZ07a}.

\begin{definition}[{\cite[Definition 2.13]{FKZ07a}}]
Через $((w_1,\ldots,w_n))$, где $w_1,\ldots,w_n\in\NN,$ мы обозначаем взвешенный граф на $n$ вершинах, который является циклом и веса вершин которого равны $w_1,\ldots,w_n$ в порядке некоторого обхода цикла. Такой взвешенный граф называется \emph{круговым}.

Круговой граф  $((w_1,\ldots,w_n))$ называется \emph{стандартным}, если выполнено одно из условий:
\begin{itemize}
    \item либо $w_1=\ldots=w_{2k}=0$ и $w_{2k+1},\ldots,w_n\leqslant-2$ для некоторого $k\geqslant0$, 
    \item либо $w_1=\ldots=w_{n-1}=0$ и $w_n\leqslant0$,
    \item либо $w_1=\ldots=w_{n-2}=0$, $w_{n-1}=w_n=-1$ и $n$ чётно.
\end{itemize}
\end{definition}

\begin{definition}[{\cite[Definitions 2.3]{FKZ07a}}]
\emph{Внутренним раздутием} в ребре $e$ называется преобразование взвешенного графа $\Gamma$, заключающееся в подразбиении ребра $e$ с помощью новой вершины веса $-1$, а также уменьшении на единицу весов вершин на концах $e$. Если ребро $e$ --- петля, то вес её конечной вершины уменьшается на 2.

\emph{Внешним раздутием} в вершине $v$ называется преобразование взвешенного графа $\Gamma$, заключающееся добавлении висячей вершины веса $-1$ и ребра, соединяющего её с $v$, а также уменьшении на единицу веса вершины $v$.

Внешним (соотв. внутренним) \emph{стягиванием} называется преобразование графа $\Gamma$, обратное к внешнему (соотв. внутреннему) раздутию.

\emph{Бирациональным преобразованием} графа $\Gamma$ называется композиция стягиваний и раздутий.
\end{definition}

\begin{remark}[{\cite[Section 3.5]{FKZ07a}}]
Пусть $(X,D)$ --- NC-пара. Раздутие на граничном дивизоре $D$ в точке пересечения двух компонент соответствует внутреннему раздутию двойственного графа $\Gamma(D)$, а в гладкой точке --- внешнему.
Бирациональные отображения NC-пары $(X,D)$, регулярные на $Y=X\setminus D$, порождают бирациональные преобразования графа $\Gamma(D)$, и наоборот (с точностью до выбора раздуваемой точки при внешнем раздутии). 
\end{remark}

\begin{definition}[{\cite[Definition 2.10]{FKZ07a}}]
Пусть вершина $v$ взвешенного графа $\Gamma$ имеет степень 2 и вес 0. При внутреннем раздутии в ребре, смежном c $v$, и последующем внутреннем стягивании вершины $v$ мы заменим вершину $v$ на новую вершину веса 0, при этом вес одной из соседних с $v$ вершин увеличится на 1, а другой --- уменьшится на 1.
Такое преобразование графа $\Gamma$ называется \emph{элементарным}.
\end{definition}

\begin{theorem}[{\cite[Corollary 3.19]{FKZ07a}}]\label{th:circular}
Любое бирациональное преобразование стандартного кругового взвешенного графа $\Gamma$ может быть представлено в виде композиции элементарных преобразований.
\end{theorem}

\section{Автоморфизмы поверхностей}
\subsection{Поверхность Маркова}

 Решения диофантова уравнения 
 \begin{equation}\label{eq:Markov}
 x^2+y^2+z^2=3xyz    
 \end{equation}
 в натуральных числах называются \emph{тройками Маркова}, а встречающиеся в них числа --- \emph{числами Маркова}. Они описаны А.А.~Марковым в 1879--1880 годах \cite{M1, Ma80}. Наиболее примечательное свойство состоит в том, что все тройки Маркова получаются из $(1,1,1)$ перестановками координат и инволюцией \[x\mapsto 3yz-x,y\mapsto y, z\mapsto z.\] Знаменитая \emph{гипотеза единственности} гласит, что \emph{каждое число Маркова является наибольшим элементом ровно в одной (упорядоченной по убыванию) тройке Маркова}.
 Различные подходы к интерпретации чисел Маркова, их контекст и история описаны в \cite{Ai13}.
 
 
Уравнение $x^2+y^2+z^2=xyz$ эквивалентно уравнению (\ref{eq:Markov}) --- его решения являются утроенными решениями  (\ref{eq:Markov}) \cite[Proposition 2.2]{Ai13}. Алгебраические гиперповерхности, задаваемые этими уравнениями в $\AA^3=\Spec\K[x,y,z]$, изоморфны; изоморфизм задаётся отображением $(x,y,z)\mapsto(3x,3y,3z)$. Любая из этих двух поверхностей называется \emph{поверхностью Маркова} \cite{ReVi17}.

Мы будем работать с поверхностями \[M_c=\{x^2+y^2+z^2-xyz=c\}\subset\AA^3,\] где $c\in\ZZ$, и назовём их \emph{поверхностями марковского типа}. Они также известны как обобщённые поверхности Маркова и как тетраэдрические поверхности Гурса. 

Также мы будем использовать двумерный алгебраический тор
\[\TT=(\K^\times)^2=\Spec\K[x^{\pm1},y^{\pm1}].\]

\subsection{Пополнения} Для двумерного алгебраического тора  мы рассмотрим пополнение \[\TT\embed\PP^1\times\PP^1=\{((x_1:x_2),(y_1:y_2))\}.\] Тор $\TT$ вкладывается в своё пополнение при помощи отображения  \[(x,y)\mapsto ((x:1),(y:1)).\]

Для поверхностей Маркова $M_c,\, c\in\K$ будем рассматривать пополнение \[M_c\subset\overline{M_c}=\{t(x^2+y^2+z^2)-xyz=ct^3\}\subset\PP^3=\{(x:y:z:t)\}.\] 





\subsection{Граничные дивизоры}
Граничный дивизор  $D_\TT=\PP^1\times\PP^1\setminus \TT$ состоит из следующих четырёх прямых, где $\infty$ обозначает $(1:0)$.
\begin{align*}
 L_1=\{x_1=0\}&=\{(0,t)\mid t\in\K\cup\{\infty\}\},\\   
 L_2=\{x_2=0\}&=\{(\infty,t)\mid t\in\K\cup\{\infty\}\},\\
 L_3=\{y_1=0\}&=\{(t,0)\mid t\in\K\cup\{\infty\}\},\\
 L_4=\{y_2=0\}&=\{(t,\infty)\mid t\in\K\cup\{\infty\}\}.
\end{align*}
  Индексы самопересечения этих прямых равны нулю. Двойственный граф дивизора --- цикл из четырёх вершин, соответствующих данным прямым, то есть цикл $((0,0,0,0))$ в обозначениях \cite[п.2.1]{FKZ07a}.
  
Граничный дивизор $\overline{D}=\overline{M_c}\setminus M_c$ состоит из трёх прямых c индексом самопересечения -1:
\begin{align*}
    L_x=\{x=t=0\},\\
    L_y=\{y=t=0\},\\
    L_z=\{z=t=0\}.
\end{align*}
При стягивании одной из них, скажем, $L_z$, получаем пополнение $\hat M_c$ с граничным дивизором $\hat D=\hat L_x\cup\hat L_y$, где $\hat L_x$ и $\hat L_y$ --- образы $L_x$ и $L_y$ соответственно. Двойственный граф $\Gamma(\hat D)$ является циклом $((0,0))$ с двумя вершинами веса 0.

\begin{remark}
Проекция $\phi\colon \PP^3\to\PP^2,(x:y:z:t)\mapsto(x:y:z),$ бирационально отображает (пополненную) поверхность Маркова $\overline{M_0}$ в $\PP^2$. Обратное отображение задаётся формулой \[\phi^{-1}:(x:y:z)\mapsto(xQ:yQ:zQ:xyz),\] где $Q=x^2+y^2+z^2$.
Коника $\{Q=0\}$ пересекает каждую из координатных прямых $x=0, y=0, z=0$ в двух точках, причём эти точки не являются точками пересечения координатных прямых.

Отображение $\phi^{-1}\colon \PP^2\to \overline{M_0}$ является композицией раздутий шести точек пересечения $Q$ с координатными прямыми и стягивания (-2)-кривой $Q$, что даёт квадратичную особую точку, то есть особенность дю Валя типа $A_1$, в $(0:0:0:1)\in\PP^3$. Образы координатных прямых $\PP^2$ образуют граничный дивизор $\overline{M_0}\setminus M_0$, и их индексы самопересечения равны -1. 
\end{remark}



\subsection{Группы автоморфизмов}
Хорошо известно (см., например, \cite[Example 1]{ArGa17}), что
\[\Aut(\TT) = \left\{(x,y)\mapsto (t_1x^ay^b,t_2x^cy^d)\mid (t_1,t_2)\in \TT, \left(\begin{smallmatrix}a&b\\c&d\end{smallmatrix}\right)\in\GL(2,\ZZ)\right\}\cong \TT\rtimes \GL(2,\ZZ).\]
Это напрямую следует из того, что автоморфизмы сохраняют подмножество обратимых функций $\{tx^ay^b\mid a,b\in\ZZ,t\in\K\}\subset \K[\TT]$. %Будем считать, что $\GL(2,\ZZ)=\{(x,y)\mapsto(x^ay^b,x^cy^d)\}$ вложена в $\Aut(\TT)$. 
Связная компонента $\Aut^\circ(\TT)=\TT$ действует тривиально на двойственном графе $\Gamma(\PP^1\times\PP^1,D_\TT)$, а дискретная часть $\Aut(\TT)/\TT$ изоморфна $\GL(2,\ZZ)$. 

Введём обозначения для некоторых элементов $\GL(2,\ZZ)$:
\begin{align*}
C_2=\begin{pmatrix}-1&0\\0&-1\end{pmatrix}&\colon(x,y)\mapsto\left(\frac{1}{x},\frac{1}{y}\right), \\    
C_4=\begin{pmatrix}0&1\\-1&0\end{pmatrix}&\colon(x,y)\mapsto\left(y,\frac{1}{x}\right), \\    
C_6=\begin{pmatrix}1&1\\-1&0\end{pmatrix}&\colon(x,y)\mapsto\left(xy,\frac{1}{x}\right). \\    
\end{align*}
Тогда $C_4^2=C_6^3=C_2$, $C_2^2=\mathrm{id}$ и, как известно \cite[Proposition 2.1]{BCCPS12}, \[\SL(2,\ZZ)=\langle C_4\rangle *_{\langle C_2 \rangle} \langle C_6\rangle\] --- амальгамированное произведение конечных циклических подгрупп, порождённых $C_4$ и $C_6$ соответственно.  




 Автоморфизм $C_4$ переставляет компоненты по циклу $L_1\to L_4\to L_2\to L_3\to L_1$, а $C_6$ является композицией пары элементарных преобразований в противоположных вершинах, а именно,
 получается раздутием точек $p_1=(0,\infty), p_2=(\infty,0)$ и стягиванием образов прямых $L_1$ и $L_2$. При раздутии точек $p_1,p_2$, чьи исключительные кривые мы обозначим через $E_1,E_2$ соответственно, получаем пополнение $T\embed  Bl_{p_1,p_2}(\PP^1\times\PP^1)$ с граничным дивизором $D_\TT'$, компоненты которого переставляются по циклу автоморфизмом $C_6$: $L_1\to E_1\to L_4\to L_2\to E_2\to L_3\to L_1$.


Введём следующие автоморфизмы аффинной поверхности $M_c$.
\begin{itemize}
    \item Замены $(x,y,z)\mapsto (\varepsilon_x x, \varepsilon_y y, \varepsilon_z z)$, где $\varepsilon_x,\varepsilon_y,\varepsilon_z\in\{-1,1\}$ и $\varepsilon_x\varepsilon_y\varepsilon_z=1$, образуют группу Клейна $K\cong(\ZZ/2\ZZ)^2$.
    \item Группа перестановок $S_3$ действует на $M_c$ перестановками координат, а её продолжение на $\overline{M_c}$ действует на граничном дивизоре перестановками компонент  $L_x, L_y, L_z$.
    \item  Отображение 
    \[\sigma\colon x\mapsto y,\; y\mapsto x,\; z\mapsto xy-z\] 
    является инволюцией на $M_c$, а её продолжение на $\overline{M_c}$ переставляет между собой прямые $L_x$ и $L_y$, прямую $L_z$ стягивает в точку, а точку пересечения $L_x$ и $L_y$ раздувает.
\end{itemize}



\begin{proposition}\label{pr:s3}
Подгруппа автоморфизмов $M_c$, регулярно продолжающихся на $\overline{M_c}$, совпадает с образом $S_3\ltimes K$ при указанном вложении.
\end{proposition}
\begin{proof}
Пусть $g\in\Aut(M_c)$ --- автоморфизм, регулярно продолжающийся на $\overline{M_c}$ и не лежащий в $S_3\ltimes K$. С точностью до умножения на элемент $S_3$ можем считать, что $g$ переводит в себя каждую из прямых $L_x, L_y, L_z$. Из этого следует, что $g^*(x)=ax+b$ для некоторых $a\in\K^\times$, $b\in\K$, и аналогично для  $g^*(y),g^*(z)$. Несложно проверить, что тогда $g^*(x)=\pm x, g^*(y)=\pm y, g^*(z)=\pm z,$ и $g\in K$. Противоречие.
\end{proof}


\begin{theorem}\label{th:aut}
Имеем $\Aut(M_c) =\langle S_3,K,\sigma\rangle\cong (\ZZ/2\ZZ)^2\rtimes \PGL(2,\ZZ)$.
\end{theorem}
\begin{proof}
Рассмотрим следующие отображения граничных дивизоров, являющиеся неразветвлёнными двойными накрытиями:
\[\psi\colon D_\TT\to \hat{D},\quad 
L_1,L_2\stackrel{\cong}{\to}\hat L_y,\quad
L_3,L_4\stackrel{\cong}{\to}\hat L_x.\]
\[\psi'\colon D_\TT'\to \overline{D},\quad 
L_1,L_2\stackrel{\cong}{\to}L_y,\quad
L_3,L_4\stackrel{\cong}{\to}L_x,\quad
E_1,E_2\stackrel{\cong}{\to}L_z.\]
При этом $\psi'$ переходит в $\psi$ при стягивании $E_1,E_2, L_z$. 

Заметим, что $C_4$ и $C_6$ индуцируют бирациональные преобразования двойственного графа $\Gamma(D_\TT)$, эквивариантные относительно $\psi$. При этом их образы порождают все элементарные преобразования двойственного графа $\Gamma(\overline{D})$. С другой стороны, действие $C_4$ на $\Gamma(D_\TT)$ совпадает с действием $\sigma$, а $C_6$ --- с действием перестановки координат $x\mapsto y\mapsto z\mapsto x$. 

Значит, по теореме~\ref{th:circular} и предложению~\ref{pr:s3}, группа $\Aut(M_c)$  порождается $S_3\ltimes K$ и $\sigma$. Поскольку $S_3$ и $\sigma$ поднимаются до преобразований $\Gamma(D_\TT)$ в себя, а $K$ действует тривиально на $\Gamma(D_\TT)$, \[\Aut(M_c)\cong K\rtimes \GL(2,\ZZ)/C_2\cong (\ZZ/2\ZZ)^2\rtimes\PGL(2,\ZZ).\]
\end{proof}

\begin{corollary}
Автоморфизмы поверхностей марковского типа сохраняют подмножество целых точек и порождаются инволюциями.
\end{corollary}

\begin{remark}
Хотя группа алгебраических автоморфизмов $\Aut(M_c)$ дискретна, голоморфная группа автоморфизмов имеет открытую орбиту \cite[Theorem 1.4]{KKL20}.
\end{remark}

\begin{example}
Рассмотрим поверхность $Y=\TT/C_2$. Её алгебра функций равна \[\K[Y]=\K[\TT]^{C_2}=\K[a,b,c]/(a^2+b^2+c^2-abc-4),\] где $a=x+\frac{1}{x},\,b=y+\frac{1}{y},\,c=xy+\frac{1}{xy}$. Следовательно, $Y=M_4$. 
\end{example}

\begin{question}
Верно ли, что множество всех точек с целыми положительными координатами на $M_4$ распадаются на орбиты вида $\Aut(M_4)\cdot(a,a,2)$, и для различных $a$ эти орбиты различны?
\end{question}



\begin{remark}
В \cite{El74} изучались поверхности, отвечающие уравнению $xyz = P(x,y,z)$, где $P(x,y,z)$ --- произвольный многочлен второй степени без свободного члена. Они называются в \cite{El74} так же поверхностями марковского типа и обладают тем же граничным дивизором из трёх прямых $L_x,L_y,L_z$.

Рассуждения из доказательства теоремы~\ref{th:aut} о преобразованиях двойственного графа применимы к данным поверхностям и влекут альтернативное доказательство \cite[Теорема 2]{El74}.
В частности, инволюции $t_1,t_2,t_3$ в \cite{El74} получаются сопряжениями инволюции $\sigma$ с помощью элементов из $S_3$.
\end{remark}



\begin{remark}
Если для взвешенного графа $((0,0,0,0))$  рассмотреть все возможные вершины, получающиеся при всевозможных цепочках внутренних раздутий, то полученное множество вершин можно естественным образом взаимно однозначно сопоставить с множеством двоично-рациональных точек на единичной окружности $[0,1]/(0\sim 1)$. А именно, исходные вершины отождествляются соответственно с $0,\frac{1}{4},\frac{1}{2},\frac{3}{4},$ а при каждом раздутии точки пересечения полученная вершина отождествляется с серединой отрезка с концами в точках, соответствующих пересекающимся кривым.

Полученное отображение эквивариантно и даёт вложение $\SL(2,\ZZ)$ в группу Томпсона $T$ кусочно-линейных двоично-рациональных преобразований единичной окружности, см. подробности \cite{GhSe87, Fo10, DiLi16}.
Аналогичное соответствие существует и для взвешенного графа $((0,0))$ и даёт вложение $\PSL(2,\ZZ)$ в $T$.
\end{remark}
